\subsection{\texorpdfstring{The star algebra \(\mathrm{L}^{\infty}(\mathrm{X},\mu)\)}{The star algebra \textbackslash mathrm\{L\}\^{}\{\textbackslash infty\}(\textbackslash mathrm\{X\},\textbackslash mu)}}

Let X be a locally compact topological space and let \(\mu\) be a positive measure on X. We consider the commutative star algebra \(\mathrm{L}^{\infty}(\mathrm{X},\mu)\) (example 4 of I, p. 103).

Lemma 6. --- Let \(g \in \mathcal{L}^{\infty}(\mathrm{X}, \mu)\). The spectrum of the class of \(g\) in \(\mathrm{L}^{\infty}(\mathrm{X}, \mu)\) is equal to the \(\mu\)-essential image of \(g\) .

Denote by \(\widetilde{g}\) the class of \(g\) in \(\mathrm{L}^{\infty}(\mathrm{X},\mu)\) and \(S\) the \(\mu\)-essential image of \(g\). Let \(z\in \mathbf{C} - \mathbf{S}\). By definition, there exists an open neighborhood \(U\) of \(z\) such that \(\mathrm{Y} = g^{-1}(\mathrm{U})\) is locally \(\mu\)-negligible. The function \(h\) defined by \(h(x) = (g(x) - z)^{-1}\) if \(x\notin \mathrm{Y}\), and \(h(x) = 0\) if \(x\in \mathrm{Y}\), belongs to \(\mathcal{L}^{\infty}(\mathrm{X},\mu)\). Its class \(\widetilde{h}\) in \(\mathrm{L}^{\infty}(\mathrm{X},\mu)\) satisfies \((\widetilde{g} -z)\widetilde{h} = 1\) since \((g(x) - z)h(x) = 1\) for all \(x\) not belonging to the locally \(\mu\)-negligible set Y. Therefore \(z\in \mathbf{C} - \mathrm{Sp}(\widetilde{g})\) .

Conversely, let \(z \in \mathbf{C} - \mathrm{Sp}(\widetilde{g})\). Let \(h \in \mathcal{L}^{\infty}(\mathrm{X},\mu)\) be a function whose class is the inverse of \(\widetilde{g} - z\) in \(\mathrm{L}^{\infty}(\mathrm{X},\mu)\) . There exists a real number \(\mathrm{M} > 0\) such that \(|h(x)| \leqslant \mathrm{M}\) locally \(\mu\)-almost everywhere, and moreover we have \((g(x) - z)h(x) = 1\) locally \(\mu\)-almost everywhere. Let U be the open ball with center \(z\) and radius \(\mathrm{M}^{-1}\) in \(\mathbf{C}\); then \(g^{-1}(\mathrm{U})\) is contained in the locally \(\mu\)-negligible set

\[
h ^ { - 1 } ( ] \mathrm{M} , + \infty [ ) \cup \{ x \in \mathrm{X} \mid ( g ( x ) - z ) h ( x ) \neq 1 \} ,
\]

hence \(z \in \mathbf{C} - \mathrm{S}\) We have proved the lemma.

PROPOSITION 4. --- Let \(g \in \mathcal{L}^{\infty}(\mathrm{X},\mu)\) . Denote by \(\widetilde{g}\) the class of \(g\) in \(\mathrm{L}^{\infty}(\mathrm{X},\mu)\) and \(\mathrm{S}\) the spectrum of \(\widetilde{g}\) . The map \(f \mapsto f \circ g\) (remark, p. 184) is the functional calculus mapping from \(\mathcal{C}(\mathrm{S})\) into \(\mathrm{L}^{\infty}(\mathrm{X},\mu)\) associated with \(\widetilde{g}\) .

Let \(f \in \mathcal{C}(\mathrm{S})\) . The function \(h = f \circ g\) is \(\mu\)-measurable (Lemma 5 of IV, p. 184) and bounded. Denote by \(\widetilde{f \circ g}\) its class in \(\mathrm{L}^{\infty}(\mathrm{X},\mu)\) . The map \(f \mapsto \widetilde{f \circ g}\) is a continuous unital involutive algebra homomorphism from \(\mathcal{C}(\mathrm{S})\) into \(\mathrm{L}^{\infty}(\mathrm{X},\mu)\) . Since \(g(x)\) belongs to \(\mathrm{S}\) locally \(\mu\)-almost everywhere (Lemma 1 of IV, p. 181 and Lemma 6), this homomorphism sends the identity map on \(\mathrm{S}\) to the class \(\widetilde{g}\) of the function \(g\) . It therefore coincides with the functional calculus mapping of \(\widetilde{g}\) (Prop. 7 of I, p. 111).
% semantic-fragments: legacy-input-seam
\relax{}
